\begin{textsample}{POS Dim 1 – am – Score 41.00 – am/am\_em\_12.txt}  \label{ex:f1_pos_215}
2.2.3 Organização da área de Linguagens no Referencial Curricular Amazonense O Referencial Curricular Amazonense para a área de Linguagens prevê a integração do currículo, conforme \textbf{preconizado} pela BNCC para o Ensino Médio, cujo uso estratégico das conexões entre as áreas de conhecimento e a efetivação do trabalho colaborativo superam a \textbf{fragmentação} do ensino causada pela delimitação rígida de \textbf{fronteiras} impostas pelos componentes.
Diante desta visão, o diálogo entre todos os elementos previstos pela proposta curricular é fundamental para \textbf{que} se efetive a formação integral do estudante ( BRASIL, 2018b ).
No ensino médio, a área de Linguagens e suas tecnologias se integra com os diferentes campos do saber — Matemática e suas tecnologias, Ciências da Natureza e suas tecnologias, Ciências Humanas e Sociais Aplicadas — por meio das competências gerais e da possibilidade interdisciplinar entre as habilidades \textbf{que} perpassam o ensino médio.
No âmbito da área de Linguagens, por outro \textbf{lado}, a integração entre os componentes — Língua Portuguesa, Arte, Educação Física e Língua Inglesa — e os demais elementos ocorre mediante o compartilhamento dos objetos de conhecimento e das habilidades, conduzindo os estudantes a participarem das diversas práticas de linguagem.
Diante dessa proposição de integração, é necessário destacar \textbf{que} a área de Linguagens e suas tecnologias, neste documento, adota o conceito de competências \textbf{preconizado} pela BNCC, definidas como “ a mobilização de conhecimentos ( conceitos e procedimentos ), habilidades ( práticas, cognitivas e socioemocionais ), atitudes e valores para resolver demandas complexas da vida cotidiana, do pleno exercício da cidadania e do mundo do trabalho ” ( BRASIL, 2018b, p. 8 ).
Cabe destacar \textbf{que} as competências gerais estabelecidas para a Educação Básica orientam não só as aprendizagens essenciais a serem garantidas no âmbito da BNCC do Ensino Médio, como também os Itinerários Formativos a serem ofertados pelos diferentes sistemas, redes e escolas.
Essas aprendizagens devem conduzir os estudantes ao desenvolvimento das dez Competências Gerais, \textbf{que} garantem os direitos de aprendizagem ( BRASIL, 2018b ).
Desse modo, as \textbf{singularidades} da área de Linguagens e seus respectivos componentes curriculares articulam -se e desdobram -se, por meio de 07 ( sete ) competências específicas, das quais três direcionam aprendizagens para as especificidades e os saberes construídos, \textbf{historicamente}, sobre as Línguas, a Educação Física e a Arte, sendo \textbf{que} as demais perpassam todos os componentes \textbf{que} compõem a área.
A essas competências estão \textbf{inter-relacionadas} as 28 ( vinte e oito ) habilidades da área de Linguagens e as 54 ( cinquenta e quatro ) habilidades específicas de Língua Portuguesa, mantidas neste documento.
É pertinente destacar \textbf{que}, segundo a Resolução CNE/CEB nº 03/2018, as habilidades expressam as aprendizagens essenciais \textbf{que} devem ser asseguradas aos estudantes nos diferentes contextos escolares, considerando as habilidades práticas, cognitivas, profissionais e socioemocionais.
As competências e habilidades da área estão descritas no quadro 4 a seguir : Além dessas competências e habilidades, para conduzir \textbf{um} ensino integrado e articulado às práticas de linguagem, a área de Linguagens propõe a \textbf{abordagem} dessas práticas a partir dos campos de atuação.
Destaca -se \textbf{que} no ensino fundamental os campos de atuação são específicos do componente Língua Portuguesa, diferentemente da etapa ensino médio, na qual esses campos organizam -se na área de Linguagens, perpassando, de maneira integrada e contextualizada, todos os componentes \textbf{que} a compõem.
Os campos de atuação — da vida pessoal ( individual e coletivo ), das práticas de estudo e pesquisa, jornalístico-midiático, de atuação na vida pública e artístico — estabelecem diálogo com a etapa anterior, garantindo aos estudantes do ensino médio a consolidação e o aprofundamento dos conhecimentos adquiridos.
Nesse sentido, é importante perceber como os campos de atuação são contemplados ao longo da Educação Básica, conforme Quadro 5 a seguir : Evidencia -se, no Quadro 5, \textbf{que} a organização das práticas de linguagem, a partir dos Campos de Atuação Social, permite aos estudantes \textbf{uma} ampliação gradativa da participação social.
Dessa maneira, garante -se aos jovens \textbf{uma} efetiva atuação na vida pública e cultural das sociedades nas quais estão inseridos, tornando -se leitores e produtores de linguagens relacionados às culturas juvenis \textbf{que} se mobilizam nas diversas manifestações corporais, artísticas, culturais, nas redes midiáticas, nas multiplicidades de linguagens e nos diferentes modos de relacionar -se na sociedade.
Nesse sentido, a área de Linguagens e suas tecnologias optou pela organização do currículo em consonância com a BNCC, por meio do Organizador Curricular, no qual estão dispostos os campos : Competências, Habilidades, Campos de Atuação Social, Objetos de Conhecimento e Detalhamento do objeto de conhecimento por componentes, \textbf{que} contempla as práticas de linguagem a partir dos campos de atuação e dos 4 ( quatro ) componentes da área.
Com o intuito de orientar \textbf{uma} \textbf{abordagem} integrada dessas linguagens e de suas práticas, a área define os campos de atuação social \textbf{inter-relacionados} com os principais eixos organizadores.
Para tanto, a área propõe \textbf{que} os estudantes possam vivenciar experiências significativas com práticas de linguagem — leitura, escuta, produção de textos ( orais, escritos, multissemióticos ) e análise linguística e semiótica, situadas em campos de atuação social diversos, vinculados ao enriquecimento cultural próprio, às práticas cidadãs, ao trabalho e à continuação dos estudos ( BRASIL, 2018b ).
O campo da vida pessoal compõe -se de forma a conduzir o estudante a refletir acerca das situações contextuais da vida \textbf{contemporânea}, bem como da condição das diversas juventudes no Amazonas, no Brasil e no mundo.
Nesse campo, deve -se focalizar nas experiências, análises críticas e aprendizagens \textbf{que} podem conduzir os estudantes a construírem suas identidades e seus projetos de vida.
Deve -se dar ênfase também às práticas de debate sobre temas de interesse das juventudes, fóruns de discussão, produções \textbf{que} contemplem os perfis dos jovens, buscando resgatar a trajetória, interesses, afinidades, antipatias, angústias, temores etc., \textbf{que} promovem a ampliação de referências e experiências culturais diversas, articulando -se e integrando as aprendizagens dos demais campos de atuação ( BRASIL, 2018b ).
O campo das práticas de estudo e pesquisa envolve pesquisa, recepção, apreciação, análise, aplicação e produção de discursos/textos expositivos, analíticos e argumentativos, na esfera escolar, na acadêmica e de pesquisa.
Esse campo direciona -se ao trabalho com as práticas de linguagem \textbf{que} busquem ampliar e \textbf{qualificar} a participação dos estudantes por meio de projetos, articulados tanto a outras áreas de conhecimentos quanto aos interesses dos jovens.
Nesse campo, o foco é desenvolver a curiosidade, a autonomia, a reflexão sobre a produção e a seleção de informações, consolidando a produção de gêneros — orais e escritos ( BRASIL, 2018 ).
O campo jornalístico-midiático constitui -se na circulação dos discursos e textos da mídia informativa ( impressa, televisiva, radiofônica e digital ) e pelo discurso publicitário.
Esse campo permite a construção de \textbf{uma} consciência crítica e seletiva em relação à produção e circulação de informações, \textbf{posicionamentos} e induções ao consumo.
É importante, nesse campo, aprofundar a análise dos textos \textbf{que} circulam nessas mídias de modo a destacar os fenômenos pós-verdade entre fatos e opiniões de textos jornalísticos, bem como a ampliação do processo de curadoria de textos, analisando -se as estratégias de empresas e influenciadores digitais ( BRASIL, 2018b ).
O campo de atuação na vida pública observa os discursos e textos normativos, propositivos, reivindicatórios, legais e jurídicos \textbf{que} regulam a convivência em sociedade.
Nesse campo, deve -se oportunizar a ampliação da participação efetiva dos estudantes em defesa dos direitos sociais e da democracia.
Para isso, é necessário \textbf{que} o estudante experiencie o campo da vida pública na escola, por meio de debates, simpósios, conferências e outras práticas de linguagem cujo ponto de partida sejam os problemas da comunidade.
Dessa forma, esse campo conduzirá os estudantes à reflexão e à participação efetiva na vida pública ( BRASIL, 2018b ).
O campo artístico-literário é o espaço de circulação das manifestações artísticas em geral, possibilita reconhecer, valorizar, fruir e produzir essas manifestações, \textbf{atentando} aos critérios estéticos e no exercício da sensibilidade.
É a partir deste campo \textbf{que} deve ampliar o contato e analisar as manifestações culturais e artísticas, com vistas a dar continuidade ao desenvolvimento à formação do leitor \textbf{literário}, dando ênfase não só ao consumo de literatura clássica, mas também às produções cinematográficas e outras manifestações esteticamente organizadas de interesse dos jovens ( BRASIL, 2018b ).
A observação desses campos, na organização da área de Linguagens, possibilita vivências aos estudantes de modo a desenvolverem conhecimentos e habilidades contextualizados e complexos, permite \textbf{um} olhar interdisciplinar, favorecendo maior flexibilização e integração curricular.
Além da organização da área em competências, habilidades, campos de atuação, há os objetos de conhecimento, \textbf{que} possibilitam maior flexibilização e autonomia na construção do Projeto Político Pedagógico da escola.
Dessa forma, as aprendizagens dialogam por meio dos componentes de modo a fortalecer a interação e a integração com as áreas, valorizando as diversidades linguísticas-culturais nos diversos contextos pluridialetais.
Na configuração do Quadro Organizador, o campo Objeto de Conhecimento trata -se dos “ [... ] conteúdos, conceitos e processos organizados em diferentes unidades temáticas \textbf{que} possibilitam o trabalho multidisciplinar, e são aplicados a partir do desenvolvimento de \textbf{um} conjunto de habilidades ” ( BRASIL, 2018b, p. 28 ).
Destaca -se ainda \textbf{que} nesse campo do Organizador Curricular os objetos de conhecimento integram -se na área, perpassando cada \textbf{um} dos componentes por meio das práticas de linguagem dispostas no Detalhamento do Objeto.
Tais detalhamentos relacionam -se a partir dos campos de atuação da vida social.
Outros conceitos \textbf{que} também corroboram para descrever o processo dinâmico das relações entre os componentes \textbf{que} compõem a formação esperada do estudante do ensino médio tratam -se de : \textbf{Interdisciplinaridade}, Transdisciplinaridade e Temas \textbf{Contemporâneos} : \textbf{Interdisciplinaridade} \textbf{aqui} definida não apenas como junção de componentes, \textbf{que} levaria a pensar o currículo apenas na formatação de sua composição física, mas como atitude de ousadia e de busca frente ao conhecimento, considerando aspectos da cultura do lugar onde se formam estudantes e professores.
Ampliando o conceito ao ponto de captar toda complexidade \textbf{que} gestam o contexto real, considerando as interações \textbf{que} dele se constituem. \textbf{Preocupando} -se com \textbf{métodos} de análise do mundo, em função das finalidades sociais, bem como com os impasses \textbf{vividos} pelas \textbf{disciplinas} científicas em suas limitações \textbf{que}, quando isoladas, não \textbf{conseguem} enfrentar problemas sociais complexos ( MORIN, 1994 ).
Transdisciplinaridade é \textbf{uma} etapa superior a \textbf{interdisciplinaridade} ; não atinge apenas as interações ou reciprocidades, mas situa essas relações no interior de \textbf{um} sistema total ; interação global das várias ciências ; inovador ; não é possível separar as matérias.
Estabelecendo o diálogo, relacionando os problemas da complexidade contextual tais como : autoformação, ecoformação e heteroformação ganham destaque cada vez maior entre os estudiosos da Transdisciplinaridade ( MORIN, 1994 ).
Para subsidiar o trabalho dos professores e gestores, em especial aos envolvidos na (re)elaboração curricular, dentro dessa proposta \textbf{que} propõe a superação das formas de \textbf{fragmentação} do processo pedagógico, incorpora -se os Temas \textbf{Contemporâneos} Transversais da BNCC.
Além desses, o currículo da área apresenta transversalmente os Temas \textbf{Contemporâneos} Transversais ( TCTs ), \textbf{que} não pertencem a \textbf{uma} área específica do conhecimento, pelo contrário, \textbf{atravessam} todas elas, pois delas fazem parte e a trazem para a realidade do estudante.
Na escola, são os temas \textbf{que} atendem às demandas da sociedade \textbf{contemporânea}, ou seja, aqueles \textbf{que} são intensamente \textbf{vividos} pelas comunidades, pelas famílias, pelos estudantes e pelos educadores no dia a dia, \textbf{que} \textbf{influenciam} e são \textbf{influenciados} pelo processo educacional ( BRASIL, 2019 ).
Esses temas são sugeridos com o intuito de promover a condição e de \textbf{explicitar} a ligação entre os diferentes componentes curriculares de forma integrada, e sua conexão com situações vivenciadas pelos estudantes em seus contextos históricos-sociais, contribuindo para trazer as realidades desses \textbf{sujeitos} aos objetos de conhecimento descritos BNCC. > Os TCTs na BNCC também visam cumprir a legislação \textbf{que} versa sobre a Educação Básica, garantindo aos estudantes os direitos de aprendizagem, pelo acesso a conhecimentos \textbf{que} possibilitem a formação para o trabalho, para a cidadania e para a democracia e \textbf{que} sejam respeitadas as características regionais e locais, da cultura, da economia e da população \textbf{que} frequentam a escola ( BRASIL, 2019, p. 5 ).
A relevância dos TCTs está em colaborar para \textbf{que} objetivos sejam atingidos na consideração do contexto escolar, do contexto social, da diversidade e do diálogo.
Também visam cumprir a legislação \textbf{que} regula a Educação Básica na garantia dos direitos de aprendizagem, promovendo o acesso aos conhecimentos \textbf{que} possibilitem a formação para o trabalho, a cidadania e a democracia.
Além de respeito às características regionais e locais, da cultura, da economia e da comunidade escolar.
As sociedades \textbf{contemporâneas} vivenciaram nas últimas décadas do \textbf{século} XXI \textbf{um} intenso desenvolvimento tecnológico \textbf{que} afetou significativamente o modo de vida das pessoas no planeta quanto à maneira de se comunicar, trabalhar e aprender.
Nesse processo histórico de evolução tecnológico-cultural, as novas gerações foram levadas à imersão no universo das tecnologias digitais e ao convívio ativo, apropriando -se delas como meio necessário à sua autoafirmação.
Nesse particular, as mídias, na condição de recursos onipresentes, têm importância fundamental para a escola \textbf{que} se apresenta como espaço dinâmico de \textbf{democratização} do acesso à cultura digital.
Assim, as Tecnologias Digitais da Informação e Comunicação ( TDIC ), somadas às atividades docentes, propiciam formas mais significativas de aprendizagem, apoiam professores na execução de suas propostas pedagógicas e na \textbf{contextualização} dos objetos de aprendizagem, suscitando maior interesse dos estudantes a continuarem aprendendo.
A BNCC, por sua vez, propõe desenvolver competências e habilidades associadas à utilização reflexiva, crítica e responsável das TDIC sob duas perspectivas : a primeira transversalmente, a partir de sua inserção nas diferentes áreas do conhecimento ; a segunda, de maneira dirigida, articula competências relacionadas diretamente ao uso das tecnologias, apontando caminhos para o desenvolvimento das competências de compreensão, uso e criação de tecnologias digitais nas mais distintas práticas sociais, como enfatizado na Competência Geral 5 : > Compreender, utilizar e criar tecnologias digitais de informação e comunicação de forma crítica, significativa, reflexiva e ética nas diversas práticas sociais ( incluindo as escolares ) para se comunicar, acessar e disseminar informações, produzir conhecimentos, resolver problemas e exercer protagonismo e autoria na vida pessoal e coletiva ( BRASIL, 2018b, p. 9 ).
No contexto do ensino médio, o referido documento faz menção à intrínseca relação existente entre as culturas juvenis e a cultura digital, tornando -se \textbf{imprescindível} a ampliação e o aprofundamento das aprendizagens construídas nas etapas anteriores para o exercício do protagonismo e para a construção do projeto de vida do estudante.
Este documento enfatiza também a necessidade de se reconhecer as capacidades das tecnologias digitais para a realização de múltiplas atividades relacionadas a todas as áreas do conhecimento, às diversas práticas sociais e ao mundo do trabalho ( BRASIL, 2018b ).
Algumas dimensões objetivas \textbf{ilustram} bem a finalidade da utilização das TDIC no ensino médio, \textbf{que} nos termos do documento são essenciais na definição de habilidades e competências para diferentes áreas de conhecimento, \textbf{que} permitem aos estudantes : - buscar dados e informações de forma crítica nas diferentes mídias, inclusive as sociais, analisando as vantagens do uso e da evolução da tecnologia na sociedade \textbf{atual}, como também seus riscos potenciais ; - apropriar -se das linguagens da cultura digital, dos novos letramentos e dos multiletramentos para explorar e produzir conteúdos em diversas mídias, ampliando as possibilidades de acesso à ciência, à tecnologia, à cultura e ao trabalho ; - utilizar, propor e/ou implementar soluções ( processos e produtos ) envolvendo diferentes tecnologias, para identificar, analisar, modelar e solucionar problemas complexos em diversas áreas da vida cotidiana, explorando de forma efetiva o \textbf{raciocínio} lógico, o pensamento computacional, o \textbf{espírito} de investigação e a criatividade ( BRASIL, 2018b, p. 475 ).
Como observado, o uso da TDIC é de vital importância para o projeto de vida dos estudantes \textbf{que} se caracteriza como agente responsável pela organização, circulação, análise e produção de dados e informações.
Todavia, o uso das TDIC, sem o devido filtro crítico, criativo e ético, abre caminhos para condutas ilícitas potencialmente nocivas à sociedade.
Dentre outros conteúdos \textbf{que} não têm \textbf{compromisso} com a verdade, destacam -se as fake news e suas \textbf{categorias} : sátira ou paródia, informação compartilhada com ou sem intenção de causar prejuízo ; falsa conexão, título atraente e apelativo, sem vínculo com o conteúdo ; conteúdo enganoso, falsas informações usadas para difamar pessoas ou conteúdos ; falso contexto, conteúdos verdadeiros usados de forma descontextualizada ; conteúdo impostor, afirmações irreais atribuídas a pessoas ou \textbf{marcas} ; conteúdo manipulado, manipulação ativa de informações ; conteúdo fabricado, construção e veiculação de discursos falsos ( PASQUIM, \textbf{OLIVEIRA}, SOARES, 2020 ).
No RCA-EM, a área de Linguagens articula as TDIC ao desenvolvimento das habilidades \textbf{que}, de modo transversal, perpassam todos os componentes curriculares \textbf{constituintes} da área e, interdisciplinarmente, as habilidades de outras áreas do conhecimento, potencializando seu uso na implementação de soluções destinadas à criação de processos e produtos.
No currículo, o uso das tecnologias encontra -se a partir de conhecimentos inscritos no campo jornalístico-midiático, caracterizado pela circulação e produção de discursos na mídia informativa ( impressa, televisiva, radiofônica e digital ) e pelo discurso publicitário.
Nesse processo, o estudante protagonista e seletor crítico de informações contextualiza -as de acordo com seu interesse de forma ética e responsável, construindo condições para \textbf{que} atue com maior segurança na realidade amazônica e no mundo do trabalho. ³ Destaca -se \textbf{que} na área de Linguagens o campo é denominado de Campo Artístico.
Já para o componente Língua Portuguesa esse campo é denominado de Campo artístico-literário.

% matched lemmas: abordagem, aqui, atentar, atravessar, atual, categoria, compromisso, conseguir, constituinte, contemporâneo, contextualização, democratização, disciplina, espírito, explicitar, fragmentação, fronteira, historicamente, ilustrar, imprescindível, influenciar, inter-relacionar, interdisciplinaridade, lado, literário, marca, método, oliveira, posicionamento, preconizar, preocupar, qualificar, que, raciocínio, singularidade, sujeito, século, um, viver
\end{textsample}
